
\vspace{1em}

\noindent
\textbf{Literature Review}

This paper contributes to the literature on how delegating grading discretion to teachers—rather than relying on standardized assessments—affects student outcomes in education markets.

The first strand of related literature examines how incentive structures and accountability systems influence grading practices. This line of research begins with \citet{hanushek1986} and is extended by \citet{jacob2003}, \citet{figlio2006}, and \citet{neal2008}, who show that when teacher evaluations depend on student performance in high-stakes exams, educators may be incentivized to concentrate their efforts on a subset of students—or worse, to inflate or manipulate scores. However, \citet{lavy2009} finds that performance pay improves both blind-graded and non-blind-graded test outcomes, providing little evidence of grade manipulation in the Israeli context. \citet{dee2019} further show that manipulation in New York’s high school exit exams declined significantly after grading was reassigned to external teachers, and the effect disappeared entirely once exams were centrally graded. Their findings highlight how policy reforms can constrain strategic grading behavior.

A second strand of literature investigates how non-standardized assessments may systematically favor certain student demographics. Building on insights from the discrimination literature \citep{becker1971,bertrand2004}, \citet{hanna2012} use a randomized control trial in India to show that teachers penalize students from lower castes, even when academic performance is held constant. More recent work highlights gender as another axis of grading bias. \citet{lavy2018} show that teachers systematically assign higher grades to girls than to equally performing boys, leading to persistent achievement gaps later in schooling. These patterns are further supported by administrative data in \citet{burgess2013}, \citet{cornwell2013}, and \citet{ferman2022}, who document large-scale grading disparities across student demographics.

A third area of research examines how grade inflation affects students’ long-term educational trajectories. \citet{betts2003} and \citet{figlio2004} document the mixed consequences of lenient grading: while inflated grades may help lower-performing students advance to the next educational stage, they can simultaneously dampen the motivation of high-achieving students by reducing the informational value of grades. More recently, \citet{denning2023} decompose the effects of grade inflation into two channels—one that facilitates progression for marginal students, and another that reduces effort among stronger students—thereby clarifying the heterogeneous implications of inflated assessments.

A central empirical challenge in this literature is distinguishing a student’s true academic ability from the grading discretion exercised by their teacher. For example, lenient teachers may attract more motivated or higher-achieving students, creating a confounding correlation between student outcomes and grading behavior. Similarly, observed score gains near high-stakes thresholds may reflect genuine effort rather than manipulation. To address these concerns, existing studies exploit institutional features that limit endogenous sorting or enable credible counterfactual comparisons. Some bias detection papers, such as \citet{lavy2018}, leverage the quasi-random assignment of students to teachers to rule out selection bias. \citet{ferman2022} exploit Brazil’s dual-grading system, in which each high-stakes exam is scored independently by both the student’s teacher and a central examiner, allowing direct measurement of grading bias. The manipulation literature often focuses on environments in which teacher-assigned grades can be benchmarked against centralized test results. \citet{jacob2003} develop statistical algorithms to detect anomalies in student answer patterns, while \citet{dee2019} compare grading patterns across centrally and locally graded exams, using a bunching analysis within a difference-in-differences framework to isolate the causal effects of accountability reforms on manipulation.


Methodologically, this paper builds on the approach of \citet{neal2008} and \cite{denning2023} by comparing grading outcomes across adjacent cohorts of students with similar observable characteristics. Specifically, I estimate how grades changed for students of comparable academic standing following the replacement of standardized exams with teacher-assigned grades. This quasi-experimental design exploits cohort-level policy variation to identify the effect of removing external grading on both the average grade level and the distribution of grades across schools. The analysis leverages a unique institutional feature of the UK education system, in which all students take standardized exams at the beginning of secondary school. Using these early scores, I construct baseline predictions of final academic performance and establish—via multiple stationarity tests—that these predictions are stable across pre-policy cohorts. This enables me to identify a school-specific degree of grade inflation, defined as the average deviation between predicted and realized grades—a quantity rarely observed in prior studies due to the absence of nationwide experiments that eliminate standardized testing.


This paper provides the first systematic evidence on between-school variation in grading patterns following a nationwide shift from standardized examinations to teacher-assigned grades. I estimate school-level grading premiums—defined as the average grade increase awarded uniformly to all students in a given school—by comparing observed outcomes with predicted performance based on pre-policy cohorts. The findings show that only about 50\% of the observed increase in grades can be explained by student-level characteristics such as gender and prior academic achievement, as emphasized in the existing literature. The remaining variation reflects substantial heterogeneity in how schools exercised grading discretion under the non-standardized regime.

In addition to documenting shifts in grading practices, this paper provides novel quantitative predictions on how the switch from standardized exams to teacher-assigned grades affects student allocation across universities. While many institutions have recently adopted test-optional admissions policies, there is surprisingly little empirical research on how alternative grading schemes influence the distribution of students across the higher education system. By modeling counterfactual placement under standardized and non-standardized regimes, I quantify how grade inflation reshapes sorting across universities, offering a new perspective on the interaction between grading policy and access to higher education.

